\documentclass[10pt,a4paper,twocolumn]{article}

\usepackage[utf8]{inputenc}
\usepackage[T1]{fontenc}
\usepackage[italian,provide=*]{babel}
\usepackage[margin=1.8cm]{geometry}
\usepackage{booktabs}
\usepackage[protrusion=true,expansion=false]{microtype}
\usepackage{enumitem}
\usepackage{xcolor}
\usepackage[colorlinks=true,linkcolor=black,urlcolor=black]{hyperref}

\setlist{nosep,leftmargin=*}
\setlength{\columnsep}{0.8cm}
\pagestyle{empty}

\title{\vspace{-1.2cm}\textbf{Progetto Finale: Pipeline di Parsing Web\\ e Valutazione con LLM-as-Judge}}

\author{
  Denis Fava \\ \small Matricola 2018687
  \and
  Riccardo Pucci \\ \small Matricola 1994172
  \and
  Alessandro Cantaressi \\ \small Matricola 2129091
}
\date{}

\begin{document}
\maketitle

\section{Obiettivo del progetto}

Il sistema realizza una pipeline completa per l'acquisizione, l'estrazione e
la valutazione di contenuto informativo da pagine web appartenenti a quattro
domini eterogenei: \texttt{en.wikipedia.org}, \texttt{www.basketball-reference.com},
\texttt{global.morningstar.com} e \texttt{it.tradingview.com}.

Data una URL, il sistema scarica la pagina, ne estrae il testo informativo
producendo una versione pulita in Markdown, e ne misura la qualità in due modi
indipendenti: con metriche automatiche calcolate rispetto a un Gold Standard
costruito a mano, e con un giudizio qualitativo prodotto da un LLM locale.
Tutto è esposto tramite API REST e tramite una Web UI, ed è orchestrato in
quattro container avviabili con un solo comando.

La scelta dei quattro domini non è casuale ed è la ragione per cui il progetto
è interessante: un'enciclopedia fatta di prosa, un archivio statistico fatto di
tabelle, un sito editoriale finanziario e una single page application che
costruisce metà pagina in JavaScript pongono problemi di parsing
strutturalmente diversi. Come si vedrà nella sezione~\ref{sec:judge}, mettono in
crisi anche le metriche di valutazione, ed è lì che il confronto fra metrica
automatica e giudizio del modello diventa istruttivo.

\section{Organizzazione del codice}

Il progetto è orchestrato da \texttt{docker-compose.yaml} e si compone di
quattro servizi: \textbf{backend} (porta 8003), \textbf{frontend} (8004),
\textbf{mariadb} (3306) e \textbf{ollama} (11434). Nessun servizio richiede
passaggi manuali dopo \texttt{docker compose up -{}-build}: lo schema del
database viene creato all'inizializzazione del container MariaDB, il modello
LLM viene scaricato dall'entrypoint del container Ollama, e il backend popola
il database leggendo i file JSON di \texttt{gs\_data/}.

\subsection*{I parser e la programmazione a oggetti}

I quattro parser condividono lo stesso ciclo di vita e sono organizzati come
gerarchia di classi. \texttt{BaseParser} è una classe astratta che fissa la
sequenza una volta sola, secondo il pattern \emph{template method}:

\begin{enumerate}
  \item \texttt{fetch} --- ottiene l'HTML, dal database in fase di valutazione
        o dalla rete altrimenti;
  \item \texttt{make\_soup} --- costruisce l'albero DOM;
  \item \texttt{prepare} --- rimuove il rumore specifico del sito;
  \item \texttt{extract\_title} --- estrae il titolo;
  \item \texttt{extract\_blocks} --- estrae i blocchi informativi;
  \item \texttt{finalize} --- compone il Markdown finale.
\end{enumerate}

Solo \texttt{extract\_blocks} è astratto: è l'unico passo davvero specifico di
ciascun sito. Le quattro sottoclassi ridefiniscono i restanti metodi soltanto
quando serve, e il vantaggio si misura in concreto: aggiungere un dominio
significa scrivere una sottoclasse e una riga nel registro, mentre la gestione
dell'HTML vuoto, la configurazione di Crawl4AI e il formato del risultato
arrivano dalla classe base, identici per tutti.

Tre ridefinizioni meritano una nota, perché documentano altrettante scoperte
fatte sul campo:

\begin{itemize}
  \item \texttt{WikipediaParser.fetch} usa una richiesta HTTP semplice invece
        di Crawl4AI. Le pagine di Wikipedia sono renderizzate lato server:
        aprire un browser headless produrrebbe lo stesso HTML impiegando
        secondi in più.
  \item \texttt{BasketballReferenceParser.prepare} riporta nel DOM le tabelle
        che il sito nasconde dentro commenti HTML e ricompone in JavaScript.
        È anche il motivo per cui il fetch restituisce \texttt{result.html} e
        non \texttt{cleaned\_html}: quest'ultimo elimina i commenti, e con essi
        metà del contenuto.
  \item \texttt{BasketballReferenceParser.parse} è l'unica sottoclasse che
        ridefinisce il template method, perché la bacheca dei premi
        (\texttt{<ul id="bling">}) va letta fra la preparazione e la pulizia,
        in un punto che la sequenza standard non prevede.
\end{itemize}

Il dispatch avviene in \texttt{parsers/registry.py}, unico punto in cui un
dominio viene associato al proprio parser: \texttt{server.py} non contiene
alcun ramo condizionale sul dominio.

\subsection*{I servizi e il livello REST}

La logica applicativa sta in \texttt{backend/src/services/}:
\texttt{database.py} (connessione e pool), \texttt{repository.py} (tutte le
query, e solo lì), \texttt{evaluator.py} (metriche), \texttt{judge.py}
(dialogo con Ollama), \texttt{bootstrap.py} (inizializzazione),
\texttt{text\_utils.py} e \texttt{markdown\_utils.py} (normalizzazione).

\texttt{server.py} resta un livello sottile di orchestrazione e validazione:
espone i quattordici endpoint richiesti, ciascuno con il proprio
\texttt{response\_model} Pydantic, e traduce gli errori dei livelli
sottostanti in codici HTTP. Il frontend, in Jinja2, dialoga esclusivamente con
le API REST e non conosce il database.

\section{Valutazione dei parser}
\label{sec:parser}

Il Gold Standard contiene \textbf{40 entry}, dieci per dominio, ciascuna con
URL, dominio, titolo, HTML grezzo e testo di riferimento copiato a mano dal
browser. La valutazione lavora sempre sull'HTML statico salvato nel database:
è ciò che la consegna richiede, ed è anche l'unico modo di ottenere numeri
confrontabili nel tempo su siti che cambiano ogni giorno.

Sono implementate due metriche. \texttt{token\_level\_eval} è quella
obbligatoria e confronta gli \emph{insiemi} di token. \texttt{sequence\_eval} è
stata aggiunta da noi e lavora sulle \emph{sequenze}: calcola precision,
recall e F1 sulla più lunga sottosequenza comune, con l'algoritmo di
\texttt{difflib.SequenceMatcher}.

La seconda esiste perché la prima ha un punto cieco preciso. Lavorando su
insiemi, ordine e ripetizioni non contano: un parser che restituisse le stesse
parole in ordine casuale, o che ripetesse tre volte lo stesso paragrafo,
otterrebbe un F1 identico a quello di un'estrazione corretta. Abbiamo
verificato il caso limite: un testo con un paragrafo ripetuto tre volte ottiene
token F1 pari a $1.000$ e sequence F1 pari a $0.737$.

\begin{table}[htbp]
\centering
\footnotesize
\setlength{\tabcolsep}{4pt}
\begin{tabular}{lccc}
\toprule
\textbf{Dominio} & \textbf{token F1} & \textbf{seq. F1} & \textbf{judge} \\
\midrule
morningstar.com          & 0.965 & 0.951 & 4.80 \\
basketball-reference.com & 0.949 & 0.922 & 2.90 \\
wikipedia.org            & 0.930 & 0.951 & 4.80 \\
tradingview.com          & 0.889 & 0.890 & 1.90 \\
\bottomrule
\end{tabular}
\caption{Metriche medie per dominio, su dieci pagine ciascuno.}
\label{tab:parser}
\end{table}

Tutti e quattro i domini superano la soglia di $0.80$ che la consegna
classifica come ``Buono''. Il dominio più difficile è TradingView, e la ragione
è strutturale: le sue schede sono griglie di dati generate da React, con nomi
di classe CSS che contengono un hash di build e cambiano a ogni rilascio. Il
parser si aggancia perciò agli attributi \texttt{data-qa-id}, che il sito usa
per i propri test automatici e che restano stabili, e adotta una whitelist
esplicita dei widget informativi anziché una blacklist del rumore.

Due osservazioni emerse durante la costruzione del Gold Standard, che
riteniamo la parte più istruttiva del lavoro.

La prima riguarda la \textbf{coerenza del criterio di copiatura}. Su
Basketball-Reference, quattro pagine su dieci erano state trascritte includendo
le tabelle statistiche stagionali, che il parser esclude di proposito, e
omettendo la coda della pagina, che il parser include. Il risultato era un F1
medio di $0.678$ con singoli valori fino a $0.096$: non un difetto del parser,
ma una divergenza di criterio fra chi aveva scritto il testo di riferimento e
chi aveva scritto il codice. Uniformato il criterio, il dominio è passato a
$0.949$.

La seconda riguarda i \textbf{caratteri invisibili}. Il testo copiato da un
browser trascina con sé marcatori Unicode di direzionalità e byte order mark,
che TradingView inserisce attorno a ogni valore numerico. Il parser li rimuove,
il testo incollato no, e poiché la metrica confronta token separati da spazio,
lo stesso numero risultava diverso da sé stesso. La normalizzazione è stata
collocata in \texttt{text\_utils.normalize\_gold\_text}, invocata nell'unico
punto in cui un \texttt{gold\_text} entra nel database --- e non dentro la
metrica: così nel database esiste una sola versione del testo, ed è quella
pulita.

\section{Valutazione dei judge}
\label{sec:judge}

Il judge dialoga con Ollama via REST su \texttt{/api/chat}, con
\texttt{"stream": false}. Le responsabilità sono separate fra i due messaggi:
il messaggio \texttt{system} contiene ruolo, criteri di penalizzazione e scala
di punteggio; il messaggio \texttt{user} contiene solo i due testi da
confrontare. Il formato è imposto a monte con \texttt{"format": "json"}, la
temperatura è bassa perché serve un giudizio ripetibile, e i testi vengono
troncati a 2000 caratteri --- consentito dalla consegna, e su CPU la differenza
fra una valutazione che termina e una che va in timeout.

Il fallback è comunque implementato, come la consegna richiede in modo
esplicito: se il modello non produce JSON leggibile o i campi attesi mancano,
il sistema restituisce punteggio minimo con \texttt{parse\_ok} a falso. La
distinzione conta, perché la percentuale di risposte valide è uno dei criteri
con cui abbiamo scelto il modello.

\begin{table}[htbp]
\centering
\footnotesize
\setlength{\tabcolsep}{4pt}
\begin{tabular}{lccc}
\toprule
\textbf{Modello} & \textbf{media} & \textbf{JSON ok} & \textbf{tempo} \\
\midrule
ministral-3:3b & 4.17 & 100\% & $\sim$13 s \\
llama3.2:3b    & 4.00 & 100\% & $\sim$8 s \\
phi4-mini      & 3.33 & 100\% & $\sim$19 s \\
qwen3:4b       & --- & 0\% & --- \\
gemma4:e2b     & --- & 0\% & --- \\
\bottomrule
\end{tabular}
\caption{Confronto fra i cinque modelli ammessi.}
\label{tab:judge}
\end{table}

Due modelli su cinque non hanno prodotto una singola risposta valida. Sono
entrambi modelli che generano token di ragionamento prima della risposta, e il
limite di generazione impostato li troncava prima che arrivassero al JSON: un
esempio concreto del fatto che ``il modello non funziona'' e ``il modello è
configurato male'' si assomigliano molto dall'esterno. Abbiamo scelto
\texttt{ministral-3:3b}, che unisce il punteggio medio più alto alla piena
affidabilità sul formato, accettando di essere più lento di
\texttt{llama3.2:3b}.

Il risultato più interessante è però nella tabella~\ref{tab:parser}, dove le
due valutazioni si contraddicono in modo sistematico. TradingView ha il
peggior F1 ($0.889$) e il peggior judge ($1.90$), ma Basketball-Reference ha
il secondo miglior F1 ($0.949$) e il secondo peggior judge ($2.90$), mentre
Wikipedia con un F1 inferiore ($0.930$) ottiene $4.80$.

La spiegazione è che le due metriche rispondono a domande diverse. La metrica
a token misura \emph{quanto} il parser ha estratto rispetto al riferimento; il
judge valuta \emph{se il risultato somiglia a un testo informativo}. Su pagine
che sono griglie di dati --- una scheda titolo, un roster di squadra --- la
seconda domanda ha una risposta strutturalmente bassa anche quando
l'estrazione è corretta, perché il modello legge un elenco di numeri e non
riconosce prosa. È esattamente il motivo per cui la consegna richiede entrambe
le valutazioni: nessuna delle due, da sola, descriverebbe il sistema.

\section{Organizzazione del database}

Il database è MariaDB, interrogato con MariaDB Connector/Python. Tutte le
query stanno in \texttt{repository.py} e usano esclusivamente segnaposto
\texttt{?} con tupla di valori: nessuna query è costruita per interpolazione
di stringa, che è la porta d'ingresso classica della SQL injection.

Lo schema, definito in \texttt{mariadb\_data/init.sql} ed eseguito
automaticamente alla creazione del container, comprende cinque tabelle.
\texttt{web\_resources} e \texttt{gold\_standard} hanno nome e colonne fissati
dalla consegna e sono legate da una relazione uno-a-uno: la stessa colonna
\texttt{url} è insieme chiave primaria e chiave esterna, con
\texttt{ON DELETE CASCADE} a implementare il requisito per cui
\texttt{DELETE /web\_resource} deve rimuovere anche il gold standard
associato. Le altre tre --- \texttt{parsed\_documents}, \texttt{evaluations},
\texttt{judgements} --- sono di nostra progettazione e servono a
\texttt{/db\_stats}, che per specifica deve leggere dati già calcolati invece
di rifare il lavoro a ogni chiamata.

Un dettaglio implementativo merita di essere riportato perché non è ovvio.
InnoDB limita una chiave a 3072 byte; in \texttt{utf8mb4} un carattere occupa
fino a quattro byte, quindi una \texttt{VARCHAR(2048)} come chiave primaria
supererebbe il limite e la \texttt{CREATE TABLE} fallirebbe. Poiché gli URL
sono ASCII per definizione (RFC 3986 impone il percent-encoding per tutto il
resto), la colonna è dichiarata \texttt{CHARACTER SET ascii COLLATE ascii\_bin}:
2048 byte, sotto il limite, con la lunghezza richiesta dalla specifica
invariata e in più un confronto case-sensitive, che per un URL è il
comportamento corretto.

Infine, una scelta di architettura dei dati che l'esperienza ci ha imposto. I
file JSON di \texttt{gs\_data/} sono la fonte di verità e il database è una
copia di lavoro, ripopolabile da zero a ogni avvio. La ragione è che i dati del
sistema non sono tutti uguali: il parsing e le metriche si possono ricalcolare
in qualunque momento, un testo di riferimento scritto a mano no. Trattare
diversamente ciò che è riproducibile da ciò che non lo è non è un dettaglio
organizzativo, ed è anche il motivo per cui la consegna chiede che i file JSON
facciano parte del progetto consegnato.

\end{document}
